% Created 2026-03-08 Sun 16:00
% Intended LaTeX compiler: xelatex
\documentclass[12pt, letterpaper]{article}
\usepackage{graphicx}
\usepackage{longtable}
\usepackage{wrapfig}
\usepackage{rotating}
\usepackage[normalem]{ulem}
\usepackage{capt-of}
\usepackage{hyperref}
\usepackage[margin=1in]{geometry}
\usepackage{amsmath}
\usepackage{amssymb}
\usepackage[T1]{fontenc}
\usepackage[utf8]{inputenc}
\usepackage{lmodern}
\usepackage{microtype}
\usepackage{booktabs}
\usepackage{array}
\usepackage{xcolor}
\usepackage{mdframed}
\usepackage{titlesec}
\usepackage{parskip}
\usepackage{hyperref}
\hypersetup{colorlinks=true, linkcolor=blue!50!black, urlcolor=blue!50!black}
\definecolor{medblue}{HTML}{2E4057}
\definecolor{accent}{HTML}{3A7CA5}
\definecolor{lightbg}{HTML}{E8F4F8}
\definecolor{substrategray}{HTML}{F5F5F5}
\definecolor{projectionblue}{HTML}{EBF4FA}
\titleformat{\section}{\large\bfseries\color{medblue}}{\thesection}{1em}{}
\titleformat{\subsection}{\normalsize\bfseries\color{accent}}{\thesubsection}{1em}{}
\newmdenv[backgroundcolor=lightbg, linecolor=accent, linewidth=1.5pt, topline=true, bottomline=true, leftline=false, rightline=false, innerleftmargin=12pt, innerrightmargin=12pt, innertopmargin=10pt, innerbottommargin=10pt]{formulabox}
\newmdenv[backgroundcolor=substrategray, linecolor=medblue, linewidth=2pt, topline=false, bottomline=false, leftline=true, rightline=false, innerleftmargin=12pt, innerrightmargin=12pt, innertopmargin=6pt, innerbottommargin=6pt]{substratebox}
\newcommand{\layerbreak}{\vspace{2em}\noindent\rule{\textwidth}{0.8pt}\vspace{2em}}
\author{Bradley K. DePuy, Independent Researcher}
\date{2026}
\title{The Medium — A Foundational Theory of Reality}
\hypersetup{
 pdfauthor={Bradley K. DePuy, Independent Researcher},
 pdftitle={The Medium — A Foundational Theory of Reality},
 pdfkeywords={},
 pdfsubject={},
 pdfcreator={Emacs 29.3 (Org mode 9.7.39)}, 
 pdflang={English}}
\begin{document}

\maketitle
\tableofcontents

\newpage

\begin{center}
\textbf{Version 0.3 (draft)  ·  DOI: \href{https://doi.org/10.5281/zenodo.18804587}{10.5281/zenodo.18804587}  ·  CC BY 4.0}
\end{center}

\begin{center}
\emph{Out of empty chaos all of creation becomes.}
\end{center}

That sentence is the theory. Everything that follows is the attempt to say it
precisely.

\begin{center}
\vspace{1em}
\begin{substratebox}
\textbf{The central clarification of v0.3:} The substrate and the observer's
framework are two distinct layers. This document is divided accordingly.
Part I describes the substrate — purely, completely, with nothing added.
Part II describes how physics emerges as a projection of that substrate
into the observer's dimensional framework. Nothing from Part II belongs
in Part I. Nothing from Part I requires Part II to be true.
\end{substratebox}
\vspace{1em}
\end{center}

\newpage
\section{PART I: THE SUBSTRATE}
\label{sec:org2854d65}
\thispagestyle{plain}
\vspace{-1em}
\noindent\textit{The substrate description is complete in itself. It contains no dimensions, no observers, no physics. It is what it is.}
\vspace{1em}
\subsection{The Foundational Claim}
\label{sec:org72e8688}

The framework rests on a single foundational claim:

\begin{quote}
The Medium always is and always was. It has no beginning and no end. It exists
outside of any temporal framework. Intrinsic distinction is an eternal property
of the Medium — not an event that occurred but a condition that always obtains.
The structure we observe is not the result of something that happened. It is the
expression of what the Medium always is.
\end{quote}

This claim does three things at once.

\textbf{It eliminates the need for a trigger.} The question of what initiated the
universe only arises if time is assumed to be primitive. If the Medium always is
and always was, there is no before and no initiation. The trigger problem
dissolves not by answering it but by identifying it as a consequence of a false
assumption.

\textbf{It eliminates time as a primitive.} The Medium exists outside of time. What we
call time is not a property of the substrate. It is downstream — a consequence
of transition, not a fundamental feature of reality.

\textbf{It reframes differentiation.} The emergence of structure is not a process that
started. It is a condition that always obtains. The ground state and the
differentiated state coexist within the Medium eternally, because intrinsic
distinction — \(\iota\) — is eternal.
\subsection{The Ground State}
\label{sec:orgfe55ac8}

The Medium at its ground state is a fully autonomous stochastic system. It
requires no inputs, no external cause, no temporal framework, and no governing
distribution. Its randomness is not randomness within a framework — it is
randomness without framework. No outcome is privileged. No state is preferred.

This is not chaos. Chaos retains structure. The ground state of the Medium
precedes structure entirely. It is meta-randomness — randomness that is random
about its own nature.

In a substrate of infinite extent with no preferred states, every possible
configuration has nonzero probability. This is not intuitive but is
mathematically unavoidable:

\begin{quote}
Every possible fluctuation — including the emergence of ordered, structured,
low-entropy configurations — is not merely possible but necessary as a feature
of the distribution itself.
\end{quote}

The universe is not something that happened against the odds. It is a necessary
expression of a substrate that has no mechanism for suppressing any possibility.
\subsection{Axioms}
\label{sec:org33173c3}

\subsubsection{Axiom I — The Eternal Medium}
\label{sec:orgde8289c}

The Medium always is and always was. It has no beginning, no end, and no
external cause.

\begin{formulabox}
\begin{equation}
\nexists\, t_0 : M \text{ begins at } t_0 \qquad \nexists\, t_1 : M \text{ ends at } t_1
\tag{1}
\end{equation}
\end{formulabox}

\emph{Status: Axiomatic.}
\subsubsection{Axiom II — Intrinsic Distinction}
\label{sec:org1eb183d}

The Medium possesses one property: intrinsic distinction, denoted \(\iota\). It
is not an operator, not a function, not a process. It is what the Medium is —
the capacity of the undifferentiated to give rise to the differentiated.

\emph{Status: Axiomatic. Nature of \(\iota\) under investigation.}
\subsubsection{Axiom III — Locality Emergence}
\label{sec:org7dbb2ab}

Location does not exist prior to differentiation. It is produced when
differentiation, propagating at the limit \(c\), gives rise to self-sustaining
structure. Before this condition is met, position and distance are genuinely
absent — not suppressed but nonexistent.

\emph{Status: Axiomatic.}
\subsection{Primitives}
\label{sec:org005009d}

\subsubsection{Primitive I — Manifestation}
\label{sec:org915dcfd}

The activation of \(\iota\) at a point in a transition is the event of
manifestation, denoted \(V\). \(V\) is either present or absent. When present,
it is identical at every occurrence — no magnitude, no internal structure, no
degrees. A quantity that is indivisible, identical at every occurrence, and
admitting no degrees has exactly one formal representation:

\begin{formulabox}
\begin{equation}
V \equiv 1
\tag{2}
\end{equation}
\centerline{\textit{Dimensionless unity. A pure number. The simplest non-trivial element of the substrate.}}
\end{formulabox}

\emph{Status: Axiomatic.}
\subsubsection{Primitive II — The Binary Field}
\label{sec:orgbb9dd3e}

With \(V \equiv 1\), the state of the Medium is completely described by a binary
field \(\phi\): at every location, manifestation either has occurred or it has
not.

\begin{formulabox}
\begin{align}
\phi &: M \to \{0, 1\} \tag{3} \\
\phi(x) &= 1 \;\text{if manifestation has occurred at}\; x, \quad 0 \;\text{otherwise} \tag{4}
\end{align}
\centerline{\textit{Both 0 and 1 are dimensionless. The substrate has no units.}}
\end{formulabox}

This is the complete formal description of the Medium's state. Nothing more is
required.

\emph{Status: Current.}
\subsubsection{Primitive III — The Manifestation Probability}
\label{sec:org33a8f6f}

Manifestation is a binary stochastic event. At any transition \(t\), a location
either manifests or it does not. The probability \(p(x)\) is a Bernoulli
parameter over the binary field — the only statistical structure the ground
state requires.

\begin{formulabox}
\begin{equation}
p(x) = \Pr(\phi(x) = 1), \qquad p : M \to [0,1]
\tag{5}
\end{equation}
\end{formulabox}

\emph{Status: Current.}
\subsubsection{Primitive IV — Transition}
\label{sec:orgf98f75c}

What physics describes as temporal evolution is, at the level of the substrate,
transition — a change from one binary field state to another. The transition
index \(t\) is a dimensionless label. Not a clock. No duration implied. No
direction implied.

\begin{formulabox}
\begin{equation}
M_t \to M_{t+1}
\tag{6}
\end{equation}
\end{formulabox}

\emph{Status: Current.}
\subsubsection{Primitive V — The Propagation Limit}
\label{sec:org2ef7325}

The substrate imposes a constraint on how rapidly manifestation can propagate
across a transition. This limit is denoted \(c\). It is dimensionless — a pure
bound on the rate of change of the binary field.

\begin{formulabox}
\begin{equation}
|\phi_{t+1}(x) - \phi_t(x)| \leq c \qquad \text{(dimensionless)}
\tag{7}
\end{equation}
\end{formulabox}

\emph{Status: Provisional.}
\subsubsection{Primitive VI — Structure Formation}
\label{sec:org2888b9c}

A self-sustaining structure forms when propagation reaches the limit exactly.
This is the threshold at which locality is born.

\begin{formulabox}
\begin{equation}
|\phi_{t+1}(x) - \phi_t(x)| = c
\tag{8}
\end{equation}
\end{formulabox}

\emph{Status: Provisional.}
\subsubsection{Primitive VII — Co-Manifestation}
\label{sec:org66b256b}

Two locations that both carry \(\phi = 1\) share a state. That is the only
substrate relationship between them. The binary field imposes no distance, no
ordering, no spatial relationship of any kind between locations that carry the
same value. Two locations may both carry \(\phi = 1\) with no chain of
\(\phi = 1\) values connecting them across the field. The field topology is not
spatial. There is no ``between'' at the substrate level.

\begin{formulabox}
\begin{equation}
x \sim y \;\iff\; \phi(x) = \phi(y) = 1
\tag{9}
\end{equation}
\centerline{\textit{Co-manifestation. No distance. No space. No ordering. Only shared state.}}
\end{formulabox}

This is prior to locality. Spatial distance between \(x\) and \(y\) is not a
substrate relationship. It is constructed downstream, after structure formation
establishes the conditions in which distance becomes meaningful.

\emph{Status: New v0.3.}
\subsection{The Entropic Arc}
\label{sec:org2363726}

The stochastic ground state and the binary field together imply a complete arc
for physical reality within the Medium.

The ground state — \(\phi = 0\) almost everywhere, fully random, no structure —
is maximum entropy by definition. The count \(W\) of arrangements consistent
with this state is maximal:

\begin{formulabox}
\begin{equation}
W_{\text{ground}} = \text{maximal}
\tag{10}
\end{equation}
\centerline{\textit{W is a dimensionless count of arrangements — a substrate quantity.}}
\end{formulabox}

The first manifestation event — \(\phi\) flipping from 0 to 1 at a single
location — is minimum local entropy by definition. One specific configuration
from a maximum entropy substrate. This resolves, without being asked, the
deepest unexplained initial condition in current physics: the universe began in
a state of extraordinarily low entropy. Within the Medium framework this is not
a brute fact. It is a necessary consequence of Axiom II.

The arc:

\begin{enumerate}
\item The Medium at maximum entropy — \(\phi = 0\) almost everywhere.
\item \(\iota\) manifests — \(V \equiv 1\) — \(\phi\) flips from 0 to 1 — minimum local entropy.
\item Structure forms at \(c\) — locality is born — patterns propagate and stabilize.
\item Complexity builds — sustained low-entropy configurations persist across transitions.
\item The trajectory continues toward maximum entropy — \(\phi\) returns toward 0.
\end{enumerate}

\begin{center}
\emph{Not a beginning and an end. A breathing.}
\end{center}

Because \(\iota\) is an eternal property of the Medium, this arc does not occur
once. It occurs wherever and whenever \(\iota\) manifests. The binary field
blinks into structure and blinks out again, across an infinite substrate that
does not notice.
\subsection{The Substrate — Complete}
\label{sec:orga3f42ec}

The substrate description is now complete. It contains:

\begin{substratebox}
\begin{itemize}
\item One substrate: the Medium
\item One property: $\iota$
\item One event: $V \equiv 1$
\item One field: $\phi \in \{0,1\}$
\item One constraint: propagation bounded by $c$
\item One relationship: co-manifestation $x \sim y$
\item One arc: ground state $\to$ structure $\to$ ground state
\end{itemize}
\vspace{0.5em}
\noindent Nothing else belongs here. Everything else is downstream.
\end{substratebox}

\layerbreak
\section{PART II: THE PROJECTION LAYER}
\label{sec:org2e9cb02}
\vspace{-1em}
\noindent\textit{Everything in Part II is downstream of the substrate. Dimensions, constants, observers, and all of physics live here — valid, precise, and complete as descriptions of projected structure. Not descriptions of the substrate itself.}
\vspace{1em}
\subsection{Physics as Projection}
\label{sec:org63c38c6}

The Medium is pre-dimensional. \(V \equiv 1\) is a pure number. The binary field
\(\phi\) is dimensionless. The propagation limit \(c\) is a dimensionless
transition bound. The transition index \(t\) is a dimensionless label.

Dimensions do not exist in the substrate. They are constructed by observers.

An observer is a self-sustaining structure in the binary field — a stable
pattern that persists across transitions, embedded in the differentiated region
it constitutes. To describe the relationships between structures like itself, it
constructs a coordinate system — a framework of units and dimensions. That
coordinate system is not a discovery about the substrate. It is a construction
of the observer.

\begin{quote}
Every dimensional quantity in physics — every constant, every observable, every
equation — is a projection of a dimensionless substrate relationship onto the
observer's dimensional coordinate system. Physics is the complete and consistent
description of those projections. It is a perfect map of the shadow. It is not
a description of the substrate that casts it.
\end{quote}
\subsubsection{The Projection Boundary}
\label{sec:org99b47da}

The boundary between the substrate and the observer's framework is the boundary
between the dimensionless and the dimensional. Everything on the substrate side
is a pure number or a binary value. Everything on the observer side carries
units. The boundary is a structural feature of the framework.

\begin{table}[htbp]
\caption{\label{tab:projection}Substrate quantities and their observer projections. Each row is a derivation target.}
\centering
\begin{tabular}{lp{7.5cm}}
\toprule
\textbf{Substrate (dimensionless)} & \textbf{Observer projection (dimensional)}\\
\midrule
\(V \equiv 1\) & \(h\) {[}action] — Planck's constant\\
\(c\) (transition bound) & \(c\) {[}length/time] — speed of light\\
\(\phi(x) \in \{0,1\}\) & quantum amplitude, probability\\
\(t\) (transition index) & time [seconds]\\
structure formation condition & spacetime geometry\\
\(W\) (arrangement count) & thermodynamic entropy [J/K]\\
propagation pattern density & energy, mass [J, kg]\\
co-manifestation \(x \sim y\) & quantum entanglement\\
\(\iota\) (intrinsic distinction) & {[}to be derived]\\
\bottomrule
\end{tabular}
\end{table}

Table \ref{tab:projection} is not a set of results. It is a research program. Each row
identifies a correspondence that must be formally derived — a demonstration that
the observer's dimensional quantity follows necessarily from the substrate
relationship on the left, given the structure of the projection boundary. Version
0.3 establishes the framework. The derivations are the work of future versions.
\subsubsection{Planck's Constant}
\label{sec:org4dacb8d}

\(V \equiv 1\) at the substrate level. Planck's constant \(h\) is the
dimensional expression of that same event, measured by an observer whose
coordinate system assigns units of action to what the substrate presents as
dimensionless unity.

\begin{formulabox}
\begin{align}
\text{Substrate:} \quad & V \equiv 1 \qquad \text{(dimensionless)} \tag{11} \\
\text{Observer:}  \quad & h = \mathrm{proj}(V) \qquad \text{([action], observer-dependent)} \tag{12}
\end{align}
\centerline{\textit{$V$ does not become $h$. $h$ is what $V \equiv 1$ looks like through a dimensional coordinate system.}}
\end{formulabox}

An open question is carried explicitly: why does the projection of \(V\) carry
units of action specifically? The answer requires a derivation of the observer's
dimensional categories from the structure of sustained binary field
configurations. That derivation is a primary target for future development.
\subsubsection{Co-Manifestation and Entanglement}
\label{sec:org1f9fdee}

At the substrate level, two locations carrying \(\phi = 1\) share a state with
no spatial relationship between them. There is no distance at the substrate
level. There is no ``between.''

When an observer embedded in differentiated structure encounters two events that
are correlated regardless of their spatial separation, it calls this
entanglement. Entanglement is the observer's projection of co-manifestation
\(x \sim y\) into a framework that expects spatial locality to govern
correlation. The apparent mystery of entanglement — correlation without causal
connection across space — dissolves at the substrate level, where space does not
yet exist and co-manifestation requires no connection at all.
\subsubsection{The General Principle}
\label{sec:org9f9908f}

The presence of dimensional constants in the equations of physics — \(h\), \(c\),
\(G\), \(k\) — is the signature of projection. Each constant marks a place where
the observer's dimensional framework has been imposed on a substrate relationship
that has none. In the substrate's own arithmetic, those constants do not appear.

The program implied by Table \ref{tab:projection}: for each row, derive the
dimensional form of the observer quantity from the substrate quantity and a
formal account of the projection boundary. If that program succeeds, every
fundamental constant of physics will have been given a substrate derivation —
and the mystery of their values resolved, not by finding deeper constants, but
by understanding that they were never fundamental to begin with.
\subsection{The Problem With Where We Started}
\label{sec:orgc21041f}

Modern physics is extraordinary. Its achievements are beyond dispute. But the
foundations were built from the wrong starting point — from the human scale
outward, not from the substrate up.

The consequences are visible at the edges. The two most precisely verified
theories in the history of science are fundamentally incompatible. Approximately
95\% of the universe by current models consists of placeholder labels for things
that do not fit. The practitioners of the field know this and say so.

The presence of dimensional constants as brute unexplained inputs is itself a
symptom. A theory built from the substrate would not require them. They appear
because physics was built from the observer's position, and the observer's
dimensional coordinate system was baked into the language from the start.
\subsection{On Existing Physics}
\label{sec:orge40c469}

The existing mathematical frameworks of physics are valid descriptions of
differentiated structure behavior within their respective domains. They are not
wrong. They are downstream. They are precise and consistent descriptions of the
observer's projection of the substrate. The Medium framework does not replace
them. It identifies the substrate from which they project.
\section{Unresolved Questions}
\label{sec:orgc7a6cfc}

These are the frontier of the framework, stated precisely and carried explicitly.

\textbf{The Nature of \(\iota\).} What is intrinsic distinction? The provisional
characterization — the property permitting local, spontaneous entropy reduction
in a fully random maximum entropy binary substrate — is the central open
question of the theory.

\textbf{The Formal Account of the Projection Boundary.} How does an observer embedded
in the binary field construct a dimensional coordinate system? What determines
which dimensional categories emerge? This is the primary task for the next stage
of formal development.

\textbf{Why Does V Project to Action?} The observer's measurement of \(V \equiv 1\)
produces a quantity with units of action. Why this dimensional combination? Until
that derivation exists, the correspondence \(V \equiv 1 \leftrightarrow h\)
{[}action] is a named target, not a result.

\textbf{The Derivation of Each Row of Table 1.} \(c\) {[}m/s] from the dimensionless
transition bound. \(G\) from structure formation geometry. Thermodynamic entropy
from the arrangement count \(W\). Each is an open derivation, explicitly
deferred.

\textbf{The Status of \(\alpha\).} The fine-structure constant \(\alpha \approx 1/137\)
is dimensionless. It cannot be a projection artifact in the same sense as \(h\)
and \(c\). Its status within the framework requires separate investigation.

\textbf{The Structure That Emerges.} What precisely is the self-sustaining binary
pattern that forms when the structure formation condition is met? The formal
definition of these patterns is the primary task of the next stage of
development.

\textbf{The Observer Boundary.} The observer is a sustained binary configuration within
\(\phi\), attempting to describe the substrate that produced her from within it.
Every formal system contains truths it cannot prove from within itself. The
Medium framework faces an analogous boundary, and acknowledges it.
\section{Summary}
\label{sec:org9d05e6f}

\begin{table}[htbp]
\caption{\label{tab:summary}Complete symbol table for version 0.3}
\centering
\begin{tabular}{llp{3.8cm}}
\toprule
\textbf{Symbol} & \textbf{Definition} & \textbf{Status}\\
\midrule
\(\iota\) & Intrinsic Distinction — one property of the Medium & Axiomatic\\
\(V \equiv 1\) & Manifestation — dimensionless unity & New v0.3\\
\(\phi : M \to \{0,1\}\) & Binary field — complete substrate description & New v0.3\\
\(x \sim y\) & Co-manifestation — shared state, no distance & New v0.3\\
\(t\) & Transition index — dimensionless label & Current\\
\(c\) & Propagation limit — dimensionless bound & Provisional\\
\(p(x)\) & Manifestation probability — Bernoulli parameter & Current\\
\(W\) & Arrangement count — dimensionless substrate entropy measure & Current\\
Projection boundary & Substrate is dimensionless; dimensions are observer constructs & New v0.3\\
\(h = \mathrm{proj}(V)\) & Planck's constant — dimensional projection of \(V \equiv 1\) & New v0.3 · Prov.\\
\bottomrule
\end{tabular}
\end{table}
\section{A Personal Note}
\label{sec:orgb3bd74f}

I am seventy-eight years old. I am not an academic. I have no credentials in
physics or mathematics. I began this work with a question that would not leave
me alone and a conviction that the question deserved a serious attempt at an
answer.

I began this work knowing I may not live to see it reach its conclusions. That
does not trouble me. What matters is that the ideas are documented, attributed,
and made available to anyone who finds them worthy of development. If this
framework contains something true it belongs to whoever can take it further. My
name attached to it is enough — for the record, and for my family.

Gregor Mendel published his work on inheritance and it sat unread for decades
before being recognized as foundational. The ideas belong to whoever finds them.
The record belongs to their origin.
\section{License and Attribution}
\label{sec:org0e5f5ed}

This work is licensed under the Creative Commons Attribution 4.0 International
License (CC BY 4.0). You are free to share and adapt this material for any
purpose, including commercially, provided you give appropriate credit to the
original author, provide a link to the license, and indicate if changes were
made.

\url{https://creativecommons.org/licenses/by/4.0/}

All versions of this document are archived on Zenodo and remain linked to the
authorship record of Bradley K. DePuy, independent researcher.

\textbf{DOI:} \href{https://doi.org/10.5281/zenodo.18804587}{10.5281/zenodo.18804587}
\end{document}
